% ============================================================
% CAPITOLO 12 — Progetto 9: Flutter e Backend: Autenticazione e CRUD
% ============================================================
\chapter{Progetto 9: Flutter e Backend: Autenticazione e CRUD}\index{Flutter!autenticazione}\index{deep link}

\section*{Cosa Costruirai}
Un'app Flutter completamente funzionante che:
\begin{itemize}
  \item Esegue login OAuth~2.0 (Google/GitHub) tramite WebView
  \item Salva i token in modo sicuro con \texttt{flutter\_secure\_storage}
  \item Mostra la lista note dall'API con CRUD completo
  \item Gestisce errori di rete e token scaduti
  \item Funziona su emulatore Android, iOS e Chrome
\end{itemize}

\textbf{Tempo stimato}: 60--90 minuti\\
\textbf{Prerequisito}: App Flutter con UI statica (Cap.~11) + Backend funzionante

% ---------------------------------------------------------
\section{Autenticazione OAuth su Mobile}

Sul web, OAuth funziona con redirect del browser. Su mobile, il flusso è diverso:

\begin{figure}[H]
\centering
\includegraphics[width=0.8\textwidth]{08-oauth-mobile}
\caption{Flusso OAuth su mobile}
\end{figure}

\begin{praticabox}[title={Pratica --- Implementa OAuth mobile}]
\begin{lstlisting}[language={}]
1. Aggiungi flutter_web_auth_2 in pubspec.yaml
2. In auth_service.dart: apri WebView con URL del backend
   OAuth, intercetta callback scheme "notesapp://",
   estrai token e salva con flutter_secure_storage
3. Nel backend: se query ?platform=mobile, redirige a
   notesapp://auth/callback?token=...&refresh=...
\end{lstlisting}
\end{praticabox}

\begin{warnbox}
Passare token come query parameter è accettabile per deep link mobili (lo scheme \texttt{notesapp://} non è intercettabile da browser). \textsc{Non} usare mai questa tecnica su URL \texttt{https://}.
\end{warnbox}

% ---------------------------------------------------------
\section{Configurazione Android e iOS}

\textbf{Android} --- aggiungi in \texttt{AndroidManifest.xml}:
\begin{lstlisting}[language=XML]
<intent-filter>
  <action android:name="android.intent.action.VIEW" />
  <category android:name="android.intent.category.DEFAULT" />
  <category android:name="android.intent.category.BROWSABLE" />
  <data android:scheme="notesapp" android:host="auth" />
</intent-filter>
\end{lstlisting}

\textbf{iOS} --- aggiungi in \texttt{Info.plist}:
\begin{lstlisting}[language=XML]
<key>CFBundleURLTypes</key>
<array>
  <dict>
    <key>CFBundleURLSchemes</key>
    <array>
      <string>notesapp</string>
    </array>
  </dict>
</array>
\end{lstlisting}

% ---------------------------------------------------------
\section{Storage Sicuro dei Token}

\begin{praticabox}[title={Pratica --- Token service}]
\begin{lstlisting}[language={}]
Crea lib/services/token_service.dart che:
1. Usa flutter_secure_storage per salvare/leggere/cancellare
2. Salva: access_token, refresh_token
3. Espone: saveTokens(), getAccessToken(),
   getRefreshToken(), clearTokens()
4. Android: EncryptedSharedPreferences
5. iOS: Keychain

Poi in api_service.dart (Dio):
- Interceptor request: legge token, aggiunge Authorization
- Interceptor response 401: prova refresh, se fallisce logout
\end{lstlisting}
\end{praticabox}

\begin{notebox}
\texttt{flutter\_secure\_storage} usa il Keychain di iOS e EncryptedSharedPreferences di Android --- gli storage crittografati nativi del sistema operativo, molto più sicuri di SharedPreferences.
\end{notebox}

% ---------------------------------------------------------
\section{Collegare le Schermate al Backend}

\begin{praticabox}[title={Pratica --- Login funzionante}]
\begin{lstlisting}[language={}]
Aggiorna login_screen.dart:
- Bottone Google -> authService.loginWithGoogle()
- Bottone GitHub -> authService.loginWithGithub()
- Durante login: CircularProgressIndicator
- Successo: naviga a /dashboard
- Errore: SnackBar con messaggio
\end{lstlisting}
\end{praticabox}

\begin{praticabox}[title={Pratica --- Dashboard con dati reali}]
\begin{lstlisting}[language={}]
Aggiorna dashboard_screen.dart:
1. Al mount, chiama notesProvider per caricare le note
2. CircularProgressIndicator durante il caricamento
3. EmptyState se non ci sono note
4. ErrorBanner se errore di rete
5. Lista NoteCard con dati reali
6. Pull-to-refresh per ricaricare
7. FAB per creare una nuova nota
\end{lstlisting}
\end{praticabox}

\begin{praticabox}[title={Pratica --- CRUD completo}]
\begin{lstlisting}[language={}]
1. note_form_screen.dart:
   - Form con titolo e contenuto
   - Validazione: titolo obbligatorio, min 3 caratteri
   - Bottone Salva: crea o aggiorna via API
   - Se editing, precompila i campi

2. note_detail_screen.dart:
   - Mostra titolo, contenuto, date
   - AppBar: azioni modifica ed elimina
   - Elimina: dialog di conferma

3. dashboard_screen.dart:
   - Tap su NoteCard -> note_detail_screen
   - Swipe-to-delete con conferma
\end{lstlisting}
\end{praticabox}

% ---------------------------------------------------------
\section{Gestione Stato con Riverpod}

L'IA ha generato i provider. Verifica questa struttura:

\begin{lstlisting}[language=Java,title={Pattern AuthNotifier con Riverpod}]
class AuthNotifier extends StateNotifier<AuthState> {
  final AuthService _authService;
  final TokenService _tokenService;

  AuthNotifier(this._authService, this._tokenService)
      : super(const AuthState()) {
    _checkAuth(); // Auto-login al riavvio
  }

  Future<void> _checkAuth() async {
    final token = await _tokenService.getAccessToken();
    if (token != null) {
      try {
        final user = await _authService.getCurrentUser();
        state = AuthState(
          status: AuthStatus.authenticated, user: user);
      } catch (_) {
        await _tokenService.clearTokens();
        state = const AuthState(
          status: AuthStatus.unauthenticated);
      }
    } else {
      state = const AuthState(
        status: AuthStatus.unauthenticated);
    }
  }
}
\end{lstlisting}

\begin{tipbox}
Verifica che \texttt{\_checkAuth()} venga chiamato nel costruttore. È il pattern ``auto-login'': se c'è un token salvato, l'app accede automaticamente senza ripetere OAuth.
\end{tipbox}

% ---------------------------------------------------------
\section{Test dell'App}

\begin{center}
\begin{tabularx}{\textwidth}{l X}
\toprule
\textbf{Funzionalità} & \textbf{Test} \\
\midrule
Login Google & Tocca $\rightarrow$ WebView $\rightarrow$ Consenso $\rightarrow$ Dashboard \\
Login GitHub & Tocca $\rightarrow$ WebView $\rightarrow$ Consenso $\rightarrow$ Dashboard \\
Auto-login & Chiudi e riapri app $\rightarrow$ ancora autenticato \\
Lista note & Le note dal backend appaiono \\
Crea nota & FAB $\rightarrow$ form $\rightarrow$ salva $\rightarrow$ nota appare \\
Modifica nota & Tocca $\rightarrow$ edit $\rightarrow$ modifica $\rightarrow$ salva \\
Elimina nota & Swipe o icona $\rightarrow$ conferma $\rightarrow$ scompare \\
Pull refresh & Tira giù $\rightarrow$ lista aggiornata \\
Errore rete & Backend spento $\rightarrow$ errore visibile \\
Token scaduto & Rinnovo automatico \\
Logout & Tocca $\rightarrow$ login $\rightarrow$ token cancellati \\
\bottomrule
\end{tabularx}
\end{center}

\begin{checkpointbox}
Se riesci a fare login, creare/modificare/eliminare una nota dall'app mobile e i dati appaiono anche nel frontend web, l'integrazione è completa.
\end{checkpointbox}

% ---------------------------------------------------------
\section*{Riepilogo}

\begin{center}
\begin{tabularx}{\textwidth}{l X}
\toprule
\textbf{Aspetto} & \textbf{Dettaglio} \\
\midrule
Auth mobile & OAuth via WebView + deep link callback \\
Token storage & flutter\_secure\_storage (Keychain/EncryptedSharedPrefs) \\
HTTP & Dio con interceptor per token e refresh \\
State & Riverpod con StateNotifier \\
CRUD & Completo: lista, crea, modifica, elimina \\
Auto-login & Token salvato $\rightarrow$ login automatico al riavvio \\
\bottomrule
\end{tabularx}
\end{center}
